\documentclass[12pt]{article}
\usepackage{amsmath, amssymb, geometry, enumitem}
\geometry{margin=1in}
\setlength{\parindent}{0pt}
\setlength{\parskip}{1em}

\title{ECON 101: Macroeconomics \vspace{5pt}\\ Solutions -- Quiz 6}
\author{Instructor: Muhebullah Karimzada}
\date{December 03,  2025}

\begin{document}
\maketitle

\section*{Part A --- Multiple Choice Questions (1 mark each)}

\setlist[enumerate]{itemsep=2pt, topsep=2pt}

\textbf{1. The Bank of Canada's primary monetary policy instrument is:}\\
\textbf{Correct answer: (C) The target overnight interest rate.}

\textbf{Explanation:}
\begin{itemize}
    \item The Bank of Canada implements monetary policy by setting a \textbf{target for the overnight interest rate} (key policy rate).
    \item Open-market operations and settlement balances are used \emph{to keep} the overnight rate near this target.
    \item The required reserve ratio is not the main instrument in Canada.
\end{itemize}

\bigskip

\textbf{2. If the Bank of Canada lowers the target overnight rate, the most immediate effect is:}\\
\textbf{Correct answer: (D) Increase in settlement balances and lower short-term rates.}

\textbf{Explanation:}
\begin{itemize}
    \item To push the overnight rate \emph{down}, the Bank of Canada typically \textbf{adds liquidity} to the system (e.g., via open-market purchases or higher settlement balances).
    \item More settlement balances increase the supply of overnight funds, lowering the overnight rate and other short-term market rates.
    \item So the immediate effect is more settlement balances and lower short-term interest rates.
\end{itemize}

\bigskip

\textbf{3. An open-market \emph{sale} of government securities by the Bank of Canada will:}\\
\textbf{Correct answer: (B) Reduce bank reserves and raise interest rates.}

\textbf{Explanation:}
\begin{itemize}
    \item When the Bank \emph{sells} government bonds, banks and dealers pay by giving up reserves/settlement balances.
    \item This \textbf{reduces the supply} of reserves in the system.
    \item With fewer reserves, the overnight rate and other interest rates tend to rise.
    \item Hence it is a contractionary move.
\end{itemize}



\textbf{4. The monetary policy transmission mechanism begins with:}\\
\textbf{Correct answer: (C) A change in the policy interest rate.}

\textbf{Explanation:}
\begin{itemize}
    \item The first step is a change in the \textbf{target overnight rate} (policy interest rate).
    \item This then affects other interest rates, the exchange rate, asset prices, and eventually spending components (C, I, NX).
    \item So the mechanism begins with the Bank's policy rate, not directly with AD or LRAS.
\end{itemize}

\bigskip

\textbf{5. If the Bank of Canada raises interest rates, the Canadian dollar tends to:}\\
\textbf{Correct answer: (B) Appreciate, reducing net exports.}

\textbf{Explanation:}
\begin{itemize}
    \item Higher Canadian interest rates make Canadian financial assets more attractive to foreign investors.
    \item Capital inflows increase demand for the Canadian dollar.
    \item The dollar \textbf{appreciates}, making exports more expensive and imports cheaper.
    \item Net exports (NX) decline.
\end{itemize}

\bigskip

\textbf{6. Which of the following is part of the Bank of Canada's inflation-control strategy?}\\
\textbf{Correct answer: (B) A 2\% inflation target midpoint.}

\textbf{Explanation:}
\begin{itemize}
    \item The Bank of Canada and the federal government agree on an inflation-control target range, with a \textbf{2\% midpoint} (typically 1\%--3\% band).
    \item The Bank does not target the money supply or exchange rate directly, and it does not impose wage ceilings.
\end{itemize}

\bigskip

\textbf{7. Contractionary monetary policy in the AD--AS model will:}\\
\textbf{Correct answer: (C) Shift AD to the left and lower real GDP in the short run.}

\textbf{Explanation:}
\begin{itemize}
    \item Higher interest rates reduce C, I, and NX, which lowers aggregate demand.
    \item In the AD--AS diagram, this appears as a \textbf{leftward shift of AD}.
    \item In the short run, with upward-sloping SRAS, real GDP falls and the price level tends to fall or rise more slowly.
\end{itemize}

\hrule
\vspace{1em}

\section*{Part B -- Problem (13 marks)}

\textbf{Problem 8 --- Taylor Rule and Transmission Mechanism}

We are given the reaction function:
\[
i = 1 + 1.2(\pi - 2) + 0.4(Y - Y^*),
\]
where:
\begin{itemize}
    \item \(i\) = target overnight rate (percent),
    \item \(\pi\) = actual inflation rate (percent),
    \item \(2\) = inflation target midpoint (percent),
    \item \(Y - Y^*\) = output gap (percent).
\end{itemize}

Data:
\[
\pi = 3.1\%, \qquad Y - Y^* = +0.8\%.
\]

\subsection*{(a) Compute the recommended Bank of Canada target overnight rate (4 marks)}

\textbf{Step 1: Compute the inflation deviation from target.}
\[
\pi - 2 = 3.1 - 2.0 = 1.1 \quad \text{(percentage points above target)}.
\]

\textbf{Step 2: Compute the inflation term \(1.2(\pi - 2)\).}
\[
1.2 \times 1.1 = 1.32.
\]

\textbf{Step 3: Compute the output gap term \(0.4(Y - Y^*)\).}
\[
Y - Y^* = +0.8 \Rightarrow 0.4 \times 0.8 = 0.32.
\]

\textbf{Step 4: Plug into the reaction function.}
\[
i = 1 + 1.2(\pi - 2) + 0.4(Y - Y^*)
= 1 + 1.32 + 0.32 = 2.64.
\]

\textbf{Answer:}
\[
\boxed{i = 2.64\% \text{ (target overnight interest rate)}}
\]

\subsection*{(b) Is the resulting policy expansionary or contractionary? Explain. (2 marks)}

\textbf{Step 1: Interpretation of the rule.}
\begin{itemize}
    \item Inflation is above target (3.1\% vs 2\%).
    \item The output gap is positive (economy operating 0.8\% above potential).
\end{itemize}

\textbf{Step 2: Policy stance.}
\begin{itemize}
    \item The rule recommends a relatively \emph{higher} interest rate (2.64\%) to cool an overheating economy.
    \item Compared to a neutral or lower rate, this implies tighter monetary conditions.
\end{itemize}

\textbf{Conclusion:}
\[
\boxed{\text{The implied stance is contractionary (tightening), intended to reduce inflationary pressure.}}
\]

\subsection*{(c) Transmission mechanism: effect of a higher policy rate (4 marks)}

\textbf{Step-by-step linkage:}

\textbf{Step 1: From policy rate to market interest rates.}
\begin{itemize}
    \item The Bank of Canada raises the target overnight rate to 2.64\%.
    \item Short-term money market rates, bank prime rates, and other lending rates rise.
    \item Longer-term interest rates may also rise as expectations adjust.
\end{itemize}

\textbf{Step 2: Effect on consumption.}
\begin{itemize}
    \item Higher borrowing costs make credit-financed purchases (cars, appliances, renovations, housing) more expensive.
    \item Some households delay or cancel spending plans.
    \item Consumption spending falls, especially on interest-sensitive goods.
\end{itemize}

\textbf{Step 3: Effect on investment.}
\begin{itemize}
    \item Firms face a higher cost of financing new capital projects.
    \item Some investment projects that were profitable at lower rates become unprofitable.
    \item Business investment in machinery, buildings, and technology falls.
\end{itemize}

\textbf{Step 4: Effect on exchange rate and net exports.}
\begin{itemize}
    \item Higher Canadian interest rates attract foreign capital (seeking higher returns).
    \item Demand for the Canadian dollar rises, causing the dollar to \textbf{appreciate}.
    \item Canadian exports become more expensive in foreign currency, and imports become cheaper for Canadians.
    \item Net exports (NX) decrease.
\end{itemize}

\textbf{Step 5: Effect on aggregate demand.}
\begin{itemize}
    \item Lower consumption (\(C\)), lower investment (\(I\)), and lower net exports (\(NX\)) all reduce total planned spending.
    \item In the AD--AS framework, this is represented by a \textbf{leftward shift in the AD curve}.
\end{itemize}

\subsection*{(d) AD--AS explanation: bringing inflation back toward 2\% (3 marks)}

\textbf{Initial situation:}
\begin{itemize}
    \item Inflation is above target (3.1\%) and the output gap is positive (\(Y > Y^*\)).
    \item In the AD--AS diagram, the economy is at an equilibrium where AD intersects SRAS to the \emph{right} of LRAS.
\end{itemize}

\textbf{Effect of contractionary monetary policy:}
\begin{itemize}
    \item Higher interest rates reduce C, I, and NX, shifting the AD curve left (from AD\(_0\) to AD\(_1\)).
    \item In the short run, this moves the equilibrium closer to the vertical LRAS line at \(Y^*\).
    \item Real GDP falls toward potential, and the price level (and inflation rate) begin to decline.
\end{itemize}

\textbf{Medium-run adjustment:}
\begin{itemize}
    \item With the output gap shrinking or closing, upward pressure on wages and prices diminishes.
    \item Inflation gradually moves back toward the 2\% target midpoint.
\end{itemize}

\textbf{Summary:}
\begin{itemize}
    \item Contractionary policy: AD shifts left.
    \item Short run: lower real GDP and lower inflation (or lower inflation growth).
    \item Medium run: economy returns to potential GDP with inflation near the 2\% target.
\end{itemize}

\hrule
\vspace{1em}

\begin{center}
\textit{End of Solutions}
\end{center}

\end{document}
